\documentclass[10pt,a4paper]{amsart}
\usepackage{amsmath,amssymb,amsthm,graphicx}
% Retained minimal preamble for the Stacks Project chapter sites-cohomology.tex (Cohomology on Sites).
% Verbatim excerpt of the chapter preamble: Xy-pic setup (native preamble lines 14-18) and the native
% theorem declarations with the mathematical macros (native preamble lines 154-217). Copied mechanically;
% nothing added, removed or redefined. The wrapper class and the amsmath/amssymb/amsthm/graphicx base
% packages are supplied by the assembler, so the class/externaldocument part of the native preamble is omitted.
% For commutative diagrams we use Xy-pic
\usepackage[all]{xy}

% We use 2cell for 2-commutative diagrams.
\xyoption{2cell}
% Theorem environments.
%
\theoremstyle{plain}
\newtheorem{theorem}[subsection]{Theorem}
\newtheorem{proposition}[subsection]{Proposition}
\newtheorem{lemma}[subsection]{Lemma}

\theoremstyle{definition}
\newtheorem{definition}[subsection]{Definition}
\newtheorem{example}[subsection]{Example}
\newtheorem{exercise}[subsection]{Exercise}
\newtheorem{situation}[subsection]{Situation}

\theoremstyle{remark}
\newtheorem{remark}[subsection]{Remark}
\newtheorem{remarks}[subsection]{Remarks}

\numberwithin{equation}{subsection}

% Macros
%
\def\lim{\mathop{\mathrm{lim}}\nolimits}
\def\colim{\mathop{\mathrm{colim}}\nolimits}
\def\Spec{\mathop{\mathrm{Spec}}}
\def\Hom{\mathop{\mathrm{Hom}}\nolimits}
\def\Ext{\mathop{\mathrm{Ext}}\nolimits}
\def\SheafHom{\mathop{\mathcal{H}\!\mathit{om}}\nolimits}
\def\SheafExt{\mathop{\mathcal{E}\!\mathit{xt}}\nolimits}
\def\Sch{\mathit{Sch}}
\def\Mor{\mathop{\mathrm{Mor}}\nolimits}
\def\Ob{\mathop{\mathrm{Ob}}\nolimits}
\def\Sh{\mathop{\mathit{Sh}}\nolimits}
\def\NL{\mathop{N\!L}\nolimits}
\def\CH{\mathop{\mathrm{CH}}\nolimits}
\def\proetale{{pro\text{-}\acute{e}tale}}
\def\etale{{\acute{e}tale}}
\def\QCoh{\mathit{QCoh}}
\def\Ker{\mathop{\mathrm{Ker}}}
\def\Im{\mathop{\mathrm{Im}}}
\def\Coker{\mathop{\mathrm{Coker}}}
\def\Coim{\mathop{\mathrm{Coim}}}

% Boxtimes
%
\DeclareMathSymbol{\boxtimes}{\mathbin}{AMSa}{"02}

%
% Macros for moduli stacks/spaces
%
\def\QCohstack{\mathcal{QC}\!\mathit{oh}}
\def\Cohstack{\mathcal{C}\!\mathit{oh}}
\def\Spacesstack{\mathcal{S}\!\mathit{paces}}
\def\Quotfunctor{\mathrm{Quot}}
\def\Hilbfunctor{\mathrm{Hilb}}
\def\Curvesstack{\mathcal{C}\!\mathit{urves}}
\def\Polarizedstack{\mathcal{P}\!\mathit{olarized}}
\def\Complexesstack{\mathcal{C}\!\mathit{omplexes}}
% \Pic is the operator that assigns to X its picard group, usage \Pic(X)
% \Picardstack_{X/B} denotes the Picard stack of X over B
% \Picardfunctor_{X/B} denotes the Picard functor of X over B
\def\Pic{\mathop{\mathrm{Pic}}\nolimits}
\def\Picardstack{\mathcal{P}\!\mathit{ic}}
\def\Picardfunctor{\mathrm{Pic}}
\def\Deformationcategory{\mathcal{D}\!\mathit{ef}}

% --- appended 2026-10-12 (Chief): the project own definitions of the reference/slogan/history
% display environments, copied mechanically from the shared Stacks preamble lines 8-12
% (\usepackage{verbatim} plus the three \newenvironment lines, which the chapter preamble keeps
% verbatim from the project preamble). Nothing added, removed or redefined: these environments
% carry no mathematical content and the project itself typesets them as comments. Without these
% three lines a counted statement whose own text uses them fails with Environment slogan undefined.
\usepackage{verbatim}
\newenvironment{reference}{\comment}{\endcomment}
\newenvironment{slogan}{\comment}{\endcomment}
\newenvironment{history}{\comment}{\endcomment}
\usepackage{cleveref}
\crefname{originalreference}{original reference}{original references}
\Crefname{originalreference}{Original reference}{Original references}
\makeatletter
\newlabel{lemma-kill-cohomology-class-on-covering}{{01E3}{source}{}{sourcepointer-1}{}}
\newlabel{lemma-kill-cohomology-class-on-covering@cref}{{[originalreference][0][]01E3}{[1][0][]source}{}{}{}}
\makeatother
\crefname{originallemma}{original lemma}{original lemmas}
\Crefname{originallemma}{Original lemma}{Original lemmas}
\makeatletter
\newlabel{sites-modules-lemma-kernel-finite-locally-constant}{{093U}{source}{}{sourcepointer-2}{}}
\newlabel{sites-modules-lemma-kernel-finite-locally-constant@cref}{{[originallemma][0][]093U}{[1][0][]source}{}{}{}}
\makeatother
\crefname{originallemma}{original lemma}{original lemmas}
\Crefname{originallemma}{Original lemma}{Original lemmas}
\makeatletter
\newlabel{sites-modules-lemma-locally-constant-finite-type}{{093N}{source}{}{sourcepointer-3}{}}
\newlabel{sites-modules-lemma-locally-constant-finite-type@cref}{{[originallemma][0][]093N}{[1][0][]source}{}{}{}}
\makeatother
\crefname{originallemma}{original lemma}{original lemmas}
\Crefname{originallemma}{Original lemma}{Original lemmas}
\makeatletter
\newlabel{sites-modules-lemma-morphism-locally-constant}{{093S}{source}{}{sourcepointer-4}{}}
\newlabel{sites-modules-lemma-morphism-locally-constant@cref}{{[originallemma][0][]093S}{[1][0][]source}{}{}{}}
\makeatother
\title{A bounded complex with locally constant finite type cohomology is locally represented by a complex of locally constant finite type sheaves}
\author{The Stacks Project Authors}
\date{stacks-project:sites-cohomology.tex:local-snapshot-20261011; GNU FDL 1.2 or later}
\begin{document}
\maketitle
\noindent\textit{Editorial note.} All delivered mathematical text is the author's own, copied from the pinned Stacks Project chapter file sites-cohomology.tex; no statement or proof step was rewritten, reordered or removed.\par\medskip
\noindent\textit{Editorial note.} The counted proof is the author's complete argument: the induction on the length of the cohomology interval, the passage to a covering making the top cohomology constant and lifting the generators of its finite module, and the induction step through the distinguished triangle, the local constancy of the kernel and the reduction of the connecting map to a map of modules.\par\medskip
\noindent\textit{Editorial note.} The label lemma-locally-constant names the counted statement, and the proof environment immediately following it in the source is delivered verbatim as its proof.\par\medskip
\noindent\textit{Editorial note.} The counted statement is carried by the author's own lemma environment in the chapter source.\par\medskip
\noindent\textit{Editorial note.} Pre-open gate: python operations/agent-authoring/tools/precard\_gate.py --source sites-cohomology.tex --label lemma-locally-constant --start 14246 --end 14313 --local-id 3575 returned RESULT=OK (L3575 free, occupancy FREE, public interval clear). Candidate taken from lane-s4/halfB-priority.json.\par\medskip
\noindent\textit{Editorial note.} The proof invokes 'a slight generalization of' a cited covering lemma with the parenthetical '(details omitted)'; the omitted verification is the passage from a single cohomology class to finitely many classes, while the mathematical content of the proof is written out, so the statement is counted and the parenthetical is recorded here verbatim.\par\medskip
\noindent\textit{Editorial note.} The statement and proof are copied from the pinned chapter file; the cited inputs are kept as external pointers with their canonical Stacks tags.\par\medskip
\noindent\textit{Allowed standard prerequisites (editorial).}\begin{itemize}
\item A locally constant finite type sheaf has a finite module of global sections. Modules on Sites, the lemma identifying a locally constant sheaf of finite type over a site with final object with the sheaf associated to a finite module. Supplies the finiteness of the module M with H\textasciicircum{}b(K) = the constant sheaf on M, and hence the choice of a finite generating set.
\item Lifting cohomology classes along a covering. Same chapter, the lemma killing a cohomology class after passing to a covering. Supplies, in the generalized form used in the proof, the covering over which the chosen generators of M lift to classes of H\textasciicircum{}b.
\item Local constancy of the kernel and of morphisms between locally constant sheaves. Modules on Sites, the lemma that the kernel of a surjection of locally constant finite type sheaves is locally constant of finite type, and the lemma that a morphism between locally constant sheaves is locally constant. Supply the local constancy of H\textasciicircum{}\{b-1\}(L) and the reduction of the map M\textasciicircum{}\{b-1\} -> Lambda\textasciicircum{}\{oplus r\} to a map of modules.
\end{itemize}\medskip
\ifdefined\hypertarget\hypertarget{sourcepointer-1}{}\fi
\noindent\textit{Source reference.} Original reference 01E3: sources/stacks-sites-cohomology/pinned-original-source.tex, source lines 715--715. Prerequisite pointer invoked at proof line(s) 14274; it names a standard fact of the same chapter that the counted argument uses, and every step of the counted proof built on it is retained verbatim in the delivered body.\par\medskip
\ifdefined\hypertarget\hypertarget{sourcepointer-2}{}\fi
\noindent\textit{Source reference.} Original lemma 093U: sources/stacks-chow/tags.txt, source lines 11527--11527. Prerequisite pointer invoked at proof line(s) 14294; it names a standard fact of a sibling chapter that the counted argument uses, and every step of the counted proof built on it is retained verbatim in the delivered body.\par\medskip
\ifdefined\hypertarget\hypertarget{sourcepointer-3}{}\fi
\noindent\textit{Source reference.} Original lemma 093N: sources/stacks-chow/tags.txt, source lines 11521--11521. Prerequisite pointer invoked at proof line(s) 14267; it names a standard fact of a sibling chapter that the counted argument uses, and every step of the counted proof built on it is retained verbatim in the delivered body.\par\medskip
\ifdefined\hypertarget\hypertarget{sourcepointer-4}{}\fi
\noindent\textit{Source reference.} Original lemma 093S: sources/stacks-chow/tags.txt, source lines 11525--11525. Prerequisite pointer invoked at proof line(s) 14303; it names a standard fact of a sibling chapter that the counted argument uses, and every step of the counted proof built on it is retained verbatim in the delivered body.\par\medskip
\begin{lemma}
\label{lemma-locally-constant}
Let $\mathcal{C}$ be a site with final object $X$.
Let $\Lambda$ be a Noetherian ring.
Let $K \in D^b(\mathcal{C}, \Lambda)$
with $H^i(K)$ locally constant sheaves of $\Lambda$-modules
of finite type. Then there exists a covering $\{U_i \to X\}$
such that each $K|_{U_i}$ is represented by
a complex of locally constant sheaves of $\Lambda$-modules
of finite type.
\end{lemma}

\begin{proof}
Let $a \leq b$ be such that $H^i(K) = 0$ for $i \not \in [a, b]$.
By induction on $b - a$ we will prove there exists a covering
$\{U_i \to X\}$ such that $K|_{U_i}$ can be represented by a complex
$\underline{M^\bullet}_{U_i}$ with $M^p$ a finite type $\Lambda$-module
and $M^p = 0$ for $p \not \in [a, b]$. If $b = a$, then
this is clear. In general, we may replace $X$ by the members
of a covering and assume that $H^b(K)$ is constant, say
$H^b(K) = \underline{M}$. By Modules on Sites, Lemma
\ref{sites-modules-lemma-locally-constant-finite-type}
the module $M$ is a finite $\Lambda$-module. Choose a surjection
$\Lambda^{\oplus r} \to M$ given by generators $x_1, \ldots, x_r$
of $M$.

\medskip\noindent
By a slight generalization of
Lemma \ref{lemma-kill-cohomology-class-on-covering} (details omitted)
there exists a covering $\{U_i \to X\}$ such that $x_i \in H^0(X, H^b(K))$
lifts to an element of $H^b(U_i, K)$. Thus, after replacing $X$ by the
$U_i$ we reach the situation where there is a map
$\underline{\Lambda^{\oplus r}}[-b] \to K$
inducing a surjection on cohomology sheaves in degree $b$.
Choose a distinguished triangle
$$
\underline{\Lambda^{\oplus r}}[-b] \to K \to L \to
\underline{\Lambda^{\oplus r}}[-b + 1]
$$
Now the cohomology sheaves of $L$ are nonzero only in the interval
$[a, b - 1]$, agree with the cohomology sheaves of $K$ in the interval
$[a, b - 2]$ and there is a short exact sequence
$$
0 \to H^{b - 1}(K) \to H^{b - 1}(L) \to
\underline{\Ker(\Lambda^{\oplus r} \to M)} \to 0
$$
in degree $b - 1$. By
Modules on Sites, Lemma
\ref{sites-modules-lemma-kernel-finite-locally-constant}
we see that $H^{b - 1}(L)$ is locally constant of finite type.
By induction hypothesis we obtain an isomorphism
$\underline{M^\bullet} \to L$ in $D(\mathcal{C}, \underline{\Lambda})$
with $M^p$ a finite $\Lambda$-module and $M^p = 0$ for
$p \not \in [a, b - 1]$. The map $L \to \Lambda^{\oplus r}[-b + 1]$
gives a map $\underline{M^{b - 1}} \to \underline{\Lambda^{\oplus r}}$
which locally is constant
(Modules on Sites, Lemma
\ref{sites-modules-lemma-morphism-locally-constant}).
Thus we may assume it is given by a map $M^{b - 1} \to \Lambda^{\oplus r}$.
The distinguished triangle shows that the composition
$M^{b - 2} \to M^{b - 1} \to \Lambda^{\oplus r}$ is zero
and the axioms of triangulated categories produce an isomorphism
$$
\underline{M^a \to \ldots \to M^{b - 1} \to \Lambda^{\oplus r}}
\longrightarrow K
$$
in $D(\mathcal{C}, \Lambda)$.
\end{proof}

\end{document}
